\chapter{Preparación de los datos GPS}
\label{ch:o1}

Este capítulo describe la primera fase del trabajo: convertir los registros GPS
en bruto del repositorio de seguimiento en una tabla regular con una posición por
individuo y día, que es la base de los modelos de los capítulos siguientes. El
criterio que guía todo el proceso es intervenir lo menos posible sobre los datos
y no añadir información que no se haya observado.

\section{Introducción y objetivo}
\label{sec:o1-intro}

Los dispositivos GPS no producen una señal continua, sino una sucesión de
\textit{fixes}: posiciones puntuales con marca temporal. Su frecuencia depende de
la programación del transmisor, de la cobertura de satélites y de la energía
disponible, de modo que la serie resultante es irregular (con intervalos que van
de unos minutos a varios días) y desigual entre individuos. Los tres modelos de
este trabajo (cadenas de Markov, modelos ocultos de Markov y aprendizaje
supervisado) comparten un requisito: operan sobre una secuencia de estados a
intervalos regulares. El objetivo de este capítulo es, por tanto, transformar esa
señal irregular y desigual en una secuencia diaria homogénea.

Dos decisiones de diseño determinan el resto del capítulo. La primera es la
\emph{granularidad diaria}: en lugar de modelar el tiempo de forma continua, cada
individuo se reduce a una posición representativa por día natural (en tiempo
universal coordinado, UTC). La irregularidad del muestreo hace inviable un modelo
de tiempo continuo, y la escala diaria capta el desplazamiento migratorio sin
depender de las lagunas del seguimiento. La segunda es la \emph{representación
explícita de los huecos}: los días sin observación válida no se eliminan, sino que
permanecen como filas con valores ausentes. De este modo, la ausencia de dato
queda registrada y no se convierte en una discontinuidad que un modelo de
transición día a día pudiera interpretar como un salto espacial.

Ambos principios se concretan en cuatro etapas encadenadas (carga, limpieza,
construcción de la tabla diaria y filtrado de individuos) que forman la
canalización de preparación de datos (figura \ref{fig:o1-pipeline}). Cada etapa es
una transformación pura sobre una tabla, y un único orquestador,
\texttt{build\_o1()}, las encadena y escribe los resultados en disco. Separar el
cálculo de la escritura permite probar cada paso de forma aislada y garantiza la
reproducibilidad del entregable.

\begin{figure}[H]
  \centering
  \includegraphics[width=\linewidth]{o1_pipeline.pdf}
  \caption[Canalización de preparación de datos]{Canalización de
  preparación de datos. A partir del fichero crudo de Movebank, las cuatro
  etapas (carga, limpieza, tabla diaria y filtrado de individuos) producen dos
  salidas: \texttt{daily.parquet}, la tabla diaria que alimenta el modelo de
  Markov, y \texttt{fixes\_clean.parquet}, los \textit{fixes} limpios que se
  reservan como material de auditoría para los objetivos posteriores.}
  \label{fig:o1-pipeline}
\end{figure}

La tabla diaria es la única entrada del modelo de Markov del capítulo siguiente,
por lo que las decisiones sobre granularidad y representación de huecos
condicionan directamente su diseño. La canalización conserva además los
\textit{fixes} limpios a su resolución original, antes del submuestreo, como
material de auditoría: de ahí se extraerán variables cinemáticas para la detección
de comportamiento y los modelos supervisados.

\section{El conjunto de datos de partida}
\label{sec:o1-dataset}

\subsection{Origen, especie y experimento}
\label{sec:o1-origen}

Los datos proceden del repositorio Movebank, del estudio \textit{Navigation
experiments in lesser black-backed gulls} \cite{movebank_gulls}, asociado a los
experimentos de navegación de Wikelski et al. \cite{wikelski2015}. Recogen el
seguimiento GPS de gaviotas sombrías (\textit{Larus fuscus}) entre 2009 y 2015.
Todos los individuos pertenecen a la misma especie: no hay variabilidad
interespecífica, de modo que las diferencias que se observen son atribuibles a la
variación individual, estacional o geográfica, no a la especie. Esta homogeneidad
simplifica la interpretación de los modelos y permite tratar a los individuos como
variantes de un mismo patrón migratorio.

\subsection{Formato y selección de variables}
\label{sec:o1-variables}

El conjunto sigue el formato estándar de Movebank, en el que cada fila es un
\textit{fix} y las columnas se agrupan en cuatro familias:

\begin{itemize}
  \item \emph{Identificación y contexto:} el identificador único del evento
  (\texttt{event-id}), el del individuo (\texttt{individual-local-identifier}) y
  el tipo de sensor (\texttt{sensor-type}), que en este conjunto es siempre GPS.
  \item \emph{Espacio-temporales:} la marca temporal (\texttt{timestamp}) y las
  coordenadas geográficas en grados decimales (\texttt{location-long} y
  \texttt{location-lat}).
  \item \emph{Calidad y control:} los indicadores \texttt{visible} y
  \texttt{manually-marked-outlier}, que señalan los registros que el propio
  repositorio marca como no visibles o como atípicos por inspección.
  \item \emph{Ambientales:} covariables de modelos climáticos externos, como la
  cobertura de vegetación baja y alta de ECMWF o la vegetación superficial de
  NCEP\@.
\end{itemize}

De ahí se selecciona el subconjunto de variables imprescindible para el modelado y
se renombra a una nomenclatura uniforme en \textit{snake\_case} (tabla
\ref{tab:o1-variables}). Se descartan los metadatos administrativos (nombre del
estudio, identificador físico de la etiqueta y nombre taxonómico, redundante al
haber una sola especie) y, en este objetivo, también las covariables ambientales:
no intervienen en la tabla diaria y, en el caso de la vegetación de NCEP, están
ausentes en todo el conjunto. El identificador de individuo se conserva como
cadena de texto (por ejemplo, \texttt{91732A}), ya que es una etiqueta y no un
valor numérico.

\begin{table}[H]
  \centering
  \caption[Selección y renombrado de variables]{Subconjunto de variables del
  formato Movebank conservadas en este capítulo, con su nombre en el proyecto y su uso. El
  resto de columnas (metadatos administrativos y covariables ambientales) se
  descarta.}
  \label{tab:o1-variables}
  \begin{tabular}{lll}
    \toprule
    \textbf{Columna Movebank} & \textbf{Nombre en el proyecto} & \textbf{Uso} \\
    \midrule
    \texttt{event-id} & \texttt{event\_id} & Identificador del \textit{fix} \\
    \texttt{timestamp} & \texttt{timestamp} & Eje temporal (UTC) \\
    \texttt{location-long} & \texttt{lon} & Longitud \\
    \texttt{location-lat} & \texttt{lat} & Latitud \\
    \texttt{individual-local-identifier} & \texttt{bird\_id} & Identificador del individuo \\
    \texttt{visible} & \texttt{visible} & Control de calidad \\
    \texttt{manually-marked-outlier} & \texttt{manually\_marked\_outlier} & Control de calidad \\
    \texttt{sensor-type} & \texttt{sensor\_type} & Tipo de sensor (siempre GPS) \\
    \bottomrule
  \end{tabular}
\end{table}

\subsection{Caracterización inicial}
\label{sec:o1-caracterizacion}

Tras la carga, el conjunto comprende 89\,867 \textit{fixes} de 126 individuos a
lo largo de algo más de seis años (tabla \ref{tab:o1-overview}). La cobertura por
individuo es marcadamente asimétrica: la mediana es de 262 \textit{fixes}, pero el
percentil 10 baja a 26 y el percentil 90 asciende a 1\,533. Coexisten, por tanto,
aves con un seguimiento denso y prolongado y aves con un rastro esporádico. Esta
disparidad anticipa la necesidad de un criterio de admisión por individuo, que se
introduce más adelante en este capítulo.

\begin{table}[H]
  \centering
  \caption[Caracterización del conjunto crudo]{Métricas agregadas del conjunto
  crudo de Movebank tras la carga y normalización.}
  \label{tab:o1-overview}
  \begin{tabular}{lr}
    \toprule
    \textbf{Métrica} & \textbf{Valor} \\
    \midrule
    \textit{Fixes} GPS totales & 89\,867 \\
    Individuos (\textit{Larus fuscus}) & 126 \\
    Mediana de \textit{fixes} por individuo & 262 \\
    Percentil 10 / 90 de \textit{fixes} por individuo & 26 / 1\,533 \\
    Primer registro (UTC) & 2009-05-25 \\
    Último registro (UTC) & 2015-08-27 \\
    \bottomrule
  \end{tabular}
\end{table}

\section{Limpieza de los \textit{fixes} GPS}
\label{sec:o1-limpieza}

La limpieza aplica primero varias comprobaciones de integridad estándar: descarta
los \textit{fixes} que el repositorio marca como no válidos (indicadores
\texttt{visible} y \texttt{manually-marked-outlier}), las coordenadas fuera del
rango geográfico físico y los registros duplicados en
\texttt{(bird\_id, timestamp)}. Ninguna de ellas elimina registros en este
conjunto, lo que indica que los datos se distribuyen ya parcialmente depurados. El
único filtro con efecto es el de velocidad, que reduce los 89\,867 \textit{fixes}
iniciales a 89\,847 (el 99,98\,\% del total).

Este filtro compara la velocidad implícita entre \textit{fixes} consecutivos del
mismo individuo con un máximo admisible. La
distancia entre dos posiciones se calcula con la fórmula de Haversine, que da la
distancia sobre la superficie de la esfera:

\begin{equation}
  d = 2R \arcsin\sqrt{\sin^2\!\frac{\Delta\phi}{2}
      + \cos\phi_1 \cos\phi_2 \sin^2\!\frac{\Delta\lambda}{2}},
  \label{eq:haversine}
\end{equation}

donde $\phi$ es la latitud, $\lambda$ la longitud (en radianes), $\Delta\phi$ y
$\Delta\lambda$ sus diferencias entre los dos \textit{fixes} y
$R = 6371{,}0088$ km el radio medio terrestre. La velocidad implícita es
$v = d/\Delta t$, con $\Delta t$ el intervalo entre ambos.

El umbral se fija en 120 km/h, valor que justifica la distribución de velocidades
(figura \ref{fig:o1-speed}). La mediana es de 0,66 km/h y el percentil 99,5 es de
50,6 km/h, en línea con la velocidad de crucero de \textit{Larus fuscus} (unos
50 km/h); solo el extremo de la cola (percentil 99,9 en 134 km/h) alcanza valores
incompatibles con el vuelo del ave. El umbral queda por encima del régimen de
crucero y de su margen con viento favorable, de modo que aísla los artefactos del
GPS sin penalizar desplazamientos legítimos.

Un único \textit{fix} erróneo contamina dos velocidades: la del segmento que llega
a él y la del que parte de él. Para no penalizar a sus vecinos legítimos, el filtro
es iterativo. En cada iteración elimina solo el \textit{fix} de mayor velocidad (el
responsable más probable) y recalcula las velocidades de los vecinos en la
siguiente, hasta que ninguna supera el umbral o hasta un tope de 20 iteraciones.
Sobre este conjunto el filtro alcanza ese tope sin converger por completo, pero su
efecto es marginal: descarta 20 \textit{fixes} (el 0,022\,\% del total).

\begin{figure}[H]
  \centering
  \includegraphics[width=0.85\linewidth]{o1_fig03_speed-distribution.png}
  \caption[Distribución de velocidades entre \textit{fixes}]{Distribución de las
  velocidades entre \textit{fixes} consecutivos del mismo individuo. Casi todos los
  valores se concentran muy por debajo de la velocidad de crucero de la especie; el
  umbral de 120 km/h (línea discontinua) deja un margen amplio sobre ese régimen y
  aísla los artefactos del GPS sin penalizar desplazamientos legítimos.}
  \label{fig:o1-speed}
\end{figure}

\section{Construcción de la tabla diaria}
\label{sec:o1-resample}

Con los \textit{fixes} ya limpios, el último paso reduce la trayectoria de cada
individuo a una posición por día natural (UTC): para cada día se toma el
\textit{fix} más cercano a una hora de referencia fija, y los días sin un
\textit{fix} suficientemente próximo quedan como huecos.

\subsection{Hora de referencia}
\label{sec:o1-hora}

La hora de referencia debe coincidir con uno de los picos de muestreo: fuera de
ellos hay muy pocos \textit{fixes}, y cualquier otra hora dejaría la mayoría de los
días sin una observación cercana. Esto reduce la elección a las cuatro ventanas en
que se concentra el muestreo (05:00, 08:00, 14:00 y 20:00 UTC), de volumen casi
idéntico (entre 15\,466 y
15\,656 \textit{fixes}; figura \ref{fig:o1-hourly}), por lo que el volumen no decide
entre ellas. La elección concreta de las 08:00 responde a un criterio biológico:
coinciden con el inicio del ciclo diario de actividad de \textit{Larus fuscus} en
buena parte de su rango migratorio, por lo que la posición a esa hora corresponde al
lugar de descanso nocturno (\textit{roost}) o al punto inmediatamente posterior al
despegue, un estado espacialmente estable. El anclaje
favorece además al modelo de Markov: observar el origen y el destino de cada
transición en el mismo punto del ciclo diario hace las transiciones comparables
entre días.

\begin{figure}[H]
  \centering
  \includegraphics[width=0.85\linewidth]{o1_fig04_hourly-coverage.png}
  \caption[Distribución horaria de los \textit{fixes}]{Número de \textit{fixes} por
  hora UTC tras la limpieza. El muestreo se concentra en cuatro ventanas (05:00,
  08:00, 14:00 y 20:00 UTC) de volumen casi idéntico; la ventana sombreada marca la
  elegida (08:00 $\pm$ 60 min).}
  \label{fig:o1-hourly}
\end{figure}

\subsection{Selección por cercanía y huecos explícitos}
\label{sec:o1-seleccion}

De cada día se conserva el \textit{fix} más cercano a la hora de referencia, medido
por su distancia temporal en minutos (\texttt{delta\_minutes}); no se interpola una
posición intermedia, en coherencia con el principio de no introducir datos no
observados. Un \textit{fix} se acepta como representante del día solo si cae dentro
de una ventana de tolerancia de $\pm 60$ min alrededor de las 08:00, holgada a
propósito para no fragmentar la serie con huecos innecesarios: con $\pm 60$ min se
conserva el 94,6\,\% de los días, frente al 66\,\% que dejaría una ventana estrecha
de $\pm 30$ min. No se amplía más porque entonces mezclaría posiciones de puntos
distintos del ciclo diario. Para
los días sin \textit{fix} válido, el procedimiento reconstruye el calendario completo de cada
individuo entre su primera y su última observación e inserta una fila por cada día
ausente, marcada como no válida y con coordenadas nulas. De este modo los huecos son
explícitos: una fila registrada, no una discontinuidad implícita.

La figura \ref{fig:o1-timeline} muestra el resultado para un individuo (91916A):
rachas largas de días válidos intercaladas con huecos puntuales o de varios días.
Esta heterogeneidad es habitual en el conjunto y es la razón de representar los
huecos en lugar de eliminarlos. Un modelo de transición día a día que encontrara dos
observaciones separadas por un hueco no anotado las trataría como días consecutivos e
introduciría un salto espacial inexistente.

\begin{figure}[H]
  \centering
  \includegraphics[width=0.85\linewidth]{o1_fig11_timeline-91916a.png}
  \caption[Línea de tiempo de un individuo]{Línea de tiempo del individuo 91916A
  tras el resample. En azul, los días con \textit{fix} válido (08:00 UTC
  $\pm$ 60 min); en rojo, los huecos. Alterna rachas largas de continuidad con
  huecos puntuales o de varios días, un patrón habitual en el conjunto.}
  \label{fig:o1-timeline}
\end{figure}

\section{Filtrado de individuos y dataset final}
\label{sec:o1-final}

\subsection{Umbral de días válidos}
\label{sec:o1-umbral}

El resample genera días válidos y huecos para cada individuo, pero no todos tienen
suficiente historia para estimar transiciones. El último filtro descarta los
individuos con menos de 30 días válidos. La figura \ref{fig:o1-validdays} muestra la
distribución de días válidos por individuo y el número de aves que supera distintos
umbrales. Con 30 días se conservan 82 de los 126 individuos (el 65\,\%). Un umbral
más laxo retendría más aves, pero a costa de admitir rastros demasiado escasos para
estimar transiciones; uno más estricto sacrificaría muestra útil. El corte en 30
días equilibra ambos extremos.

\begin{figure}[H]
  \centering
  \includegraphics[width=0.85\linewidth]{o1_fig06_valid-days-per-bird.png}
  \caption[Días válidos por individuo]{Distribución del número de días válidos por
  individuo y número de aves conservadas según el umbral. El corte en 30 días
  descarta los individuos con seguimiento insuficiente para el modelo de Markov sin
  penalizar a la mayoría del conjunto.}
  \label{fig:o1-validdays}
\end{figure}

\subsection{El conjunto de datos final}
\label{sec:o1-dataset-final}

La tabla \ref{tab:o1-funnel} resume la reducción de los datos a lo largo de la
pipeline.

\begin{table}[H]
  \centering
  \caption[Reducción de los datos a lo largo de la pipeline]{Reducción del volumen
  de datos en cada etapa, desde los \textit{fixes} en bruto hasta la tabla diaria
  final.}
  \label{tab:o1-funnel}
  \begin{tabular}{lr}
    \toprule
    \textbf{Etapa} & \textbf{Recuento} \\
    \midrule
    \textit{Fixes} GPS brutos (126 individuos) & 89\,867 \\
    \textit{Fixes} tras la limpieza & 89\,847 \\
    Individuos conservados ($\geq 30$ días válidos) & 82 \\
    Filas de la tabla diaria final & 24\,444 \\
    Filas válidas (con posición real) & 21\,823 \\
    \bottomrule
  \end{tabular}
\end{table}

La pipeline produce dos ficheros: \texttt{daily.parquet}, la tabla diaria que
alimenta el modelo de Markov, y \texttt{fixes\_clean.parquet}, los \textit{fixes}
limpios a su resolución original, reservados como material de auditoría para los
objetivos posteriores. Ambos se regeneran con una sola llamada a
\texttt{build\_o1()}; sus funciones de carga, limpieza y resample son puras y solo
el orquestador escribe a disco.

\section{Caracterización del dataset resultante}
\label{sec:o1-resultante}

La tabla diaria es la entrada del modelo de Markov del capítulo siguiente. Tres
propiedades del conjunto resultante condicionan cómo puede explotarse: su cobertura
espacial, la fragmentación de las series y el sesgo estacional del seguimiento.

\subsection{Dominio espacial}
\label{sec:o1-espacial}

Los \textit{fixes} conservados cubren el corredor migratorio completo de la especie,
desde las zonas de cría en el norte de Europa hasta las áreas de invernada en el
este de África (figura \ref{fig:o1-spatial}). Esta extensión fija el dominio que
la cadena de Markov deberá discretizar: la rejilla de celdas ha de abarcar desde las
latitudes nórdicas hasta el este de África.

\begin{figure}[H]
  \centering
  \includegraphics[width=0.85\linewidth]{o1_fig08_spatial-overview.png}
  \caption[Distribución geográfica de los \textit{fixes}]{Distribución geográfica de
  una muestra de los \textit{fixes} conservados. El conjunto cubre el corredor
  migratorio de \textit{Larus fuscus} entre el norte de Europa y el este de
  África.}
  \label{fig:o1-spatial}
\end{figure}

\subsection{Fragmentación de las series}
\label{sec:o1-fragmentacion}

Las rachas de días válidos consecutivos son en su mayoría cortas: la mediana es de
8 días (figura \ref{fig:o1-streak}). Para el modelo de Markov esto supone abundancia
de transiciones de un solo paso, pero escasez de tramos largos continuos en la mayor
parte de la muestra. Queda abierta para ese modelo una decisión: entrenar sobre
todas las rachas o exigir una longitud mínima.

\begin{figure}[H]
  \centering
  \includegraphics[width=0.85\linewidth]{o1_fig09_streak-length-distribution.png}
  \caption[Longitud de las rachas de días válidos]{Distribución de la longitud de
  las rachas de días válidos consecutivos en la tabla diaria (eje recortado a 200
  días). La mayoría son cortas (mediana de 8 días), con una cola larga de aves
  seguidas de forma casi continua durante meses.}
  \label{fig:o1-streak}
\end{figure}

\subsection{Sesgo estacional del seguimiento}
\label{sec:o1-estacional}

El seguimiento no se reparte por igual a lo largo del año (figura
\ref{fig:o1-seasonal}). El otoño reúne casi el doble de observaciones válidas que la
primavera, con un máximo en septiembre y un mínimo en mayo; en términos de
fenología, la migración de otoño queda mucho mejor muestreada que la de primavera.
En consecuencia, las transiciones que el modelo estime para las estaciones bien
representadas se apoyarán en más datos que las de las estaciones escasas, algo a
tener en cuenta al interpretar su precisión por régimen.

\begin{figure}[H]
  \centering
  \includegraphics[width=0.85\linewidth]{o1_fig10_monthly-seasonal-coverage.png}
  \caption[Cobertura mensual y estacional]{Filas diarias válidas por mes y por
  estación (hemisferio norte). El seguimiento se concentra en el otoño y es escaso
  en primavera, un sesgo que condiciona la interpretación de las transiciones por
  estación.}
  \label{fig:o1-seasonal}
\end{figure}
